
Here, we collect results for different classes of constraint languages. For some classes
we classify the computational complexity of the constrained synchronization problem.

\begin{enumerate}
	\item In Section~\ref{sec:constr_sync:infinite} we show that for infinite languages the problem is~\LOGSPACE-hard. 
	
	\item \Cref{sec:constr_sync:letter_bounded} investigates
	sparse languages. We show that the problem is in \NP\  for sparse constraint languages,
    in $\AC^0$ for the subclass of finite languages,
    and we state sufficient conditions for containment in \LOGSPACE.	
	we conclude with a full classification for the subclass of letter-bounded constraint languages, characterizing
	precisely when the problem is in $\AC^0$, \LOGSPACE-complete or \NP-complete.
		
	\item In \SEC~\ref{constr_sync:sec:monoid}
	we investigate languages that are submonoids of $\Sigma^*$.
	We show that algorithmic results and bounds on shortest synchronizing words
	from the unconstrained setting can be transferred to this setting.
	In particular, for these constrained problems the problem is 
	always in \NL\ (characterizing \NL-hardness in this case is an open problem).
	From this result, we also deduce in \cref{sec:constr_sync:SCC}
	that constraint languages accepted by strongly connected PDFAs
	yield constrained synchronization problems in \NL.
		 
	\item \Cref{sec:constr_sync:small_aut} is concerned with small constraint PDFAs.
	\Cref{constr_sync:sec:two-state} gives a complexity classification
	for two-state PDFAs that are not strongly connected.	
	In \cref{sec:constr_sync:three_states}, for three-state binary constraint PDFAs 
	that are not strongly connected we give a 
	full classification into problems in $\AC^0$ and problems that are \LOGSPACE-, \NL-,
	\NP-, or \PSPACE-complete.	
	
	\item Lastly, Section~\ref{sec:constr_sync:commutative} is concerned with commutative constraint languages.
	We give a complete classification, stating precisely the conditions for containment in $\AC^0$,
	\LOGSPACE-completeness, \NL-completeness, \NP-completeness,	or \PSPACE-completeness.	
\end{enumerate}

The complexity classifications for each class we give can all be used to give decision procedures
to decide the complexity for a given constraint language from the class.
